\documentclass[14pt]{extarticle}
\usepackage[margin=1in]{geometry}
\usepackage{hyperref,xcolor}
\definecolor{crimson}{rgb}{0.86, 0.08, 0.24}
\hypersetup{
    colorlinks=true,
    linkcolor=crimson,
    filecolor=crimson,      
    urlcolor=crimson,
    citecolor=crimson,
}
\urlstyle{same}

\begin{document}
Our quantity of interest is the acoustic $Q$ factor of a pipe\footnote{From \href{https://en.wikipedia.org/wiki/Q_factor\#Acoustical_systems}{wikipedia}, ``The $Q$ factor is generally a dimensionless parameter that describes how underdamped an oscillator or resonator is''}, and this value is calculated from finding three frequencies associated with the pipe:
\begin{equation}
  Q = \frac{f_0}{f_2 - f_1}
\end{equation}
where $f_0$ is some resonant frequency, $f_2 > f_1$, and where $A(f)$ represents the sound amplitude for a given frequency, $A(f_0) = A(f_1) \sqrt{2} = A(f_2) \sqrt{2}$. Thus, this value can be measured for a pipe by obtaining an amplitude vs. frequency curve for a given pipe, fitting this curve based on some analytical function similar to
\begin{equation}
  A(f) = \frac{A_0}{\sqrt{(f_0^2 - f^2)^2 + (2 \gamma f)^2}}
\end{equation}
where $A_0$ is the resonant amplitude (the maximal amplitude), $f_0$ is the resonant frequency, and $\gamma$ is some damping term\footnote{This is also from \href{https://en.wikipedia.org/wiki/Resonance\#Acoustic}{wikipedia}.}. This fit can use the initial conditions of direct observations, such as immediately guessing the resonant frequency from the length of the pipe using the open-open standing wave equation
\begin{equation}
  f_0 = \frac{v}{2L}
\end{equation}
or a guess of the resonant frequency can be made from the data. The maximal amplitude would have to be guessed from the data, and the damping term would likely need to be guessed entirely using literature and the material of the pipe. From this fit, we can extract the resonant frequency alongside $f_1$ and $f_2$, yielding $Q$ for the pipe.

Thus, our independent quantity is obviously the input frequency, and the dependent quantity is the amplitude, as for right now, we are measuring the frequency-amplitude relationship to calculate $Q$. In future, we may extend this experiment to more pipes if this data collection process can be streamlined, and then it would be seeing how pipe characteristics alter the $Q$ score.

In yesterday's lab, we were only doing familiarization. We manually sampled a  frequency-amplitude curve for a single pipe to get used to the measurement process. This manual sampling took far too long to comfort, taking the entire lab time to complete, so we were unable to perform a fit or calculate a $Q$ value. However, we learned quite a bit about how we can reduce uncertainty, such as measuring the frequency from the oscilloscope rather than the from the function generator. For next week, I hope that we can measure a full curve, fit it, and calculate $Q$ for our current pipe. If we have extra time as well, I'd also like to repeat this process for another pipe and then compare the relationship of $Q$ values across both pipes to literature expectations.

Lastly, our planned improvements for next week will need to be in uncertainty/ noise-reduction techniques. First and foremost, the room is far too noisy. I hope that bringing wool or wrapping the apparatus in clothing might be enough to promote noise isolation, but this is unlikely. We also need to have some method for sampling the frequency-amplitude resonance data to a fairly high granularity to accurately determine the resonant frequency (as this value is dependent on the speed of sound in the room, which changes significantly depending on the temperature, as we saw in our most recent experiment).
\end{document}
